\section{Conclusion}
Meeting-risk calibration is subject to two different constraints: what prices distinguish and what a model permits. In the Gaussian model, post-cutoff cash claims retain only total earlier variance, even with continuous background risk and American exercise. Later-accruing futures retain its first time moment as well. Earlier expiries can separate more of the allocation, according to the rank of their meeting and background exposures; finite premium precision determines whether that separation is useful.

For scheduled-rate factor models, arbitrary deterministic scale recalibration leaves a finite variance cone. Under constant ordinal loadings, its complete restrictions are consecutive-meeting inequalities. With square-root stochastic volatility, the full drift-and-leverage kernel gives an affine variance map whose exact support and covariance rank depend on the active variance states. The two published three-factor specifications consequently impose different numbers of affine restrictions. These conclusions concern model restrictions, independently of whether a particular option panel identifies the underlying meeting variances.

The settlement application shows why both distinctions belong in a calibration report. A smile fit sharply reduces pricing errors, yet changing the background still removes positive meeting-risk lower bounds. A known nonparallel loading lets same-expiry options on different futures add information; estimating that loading creates another source of uncertainty. Allowing a free long-end loading improves three-year-midcurve predictions relative to a hump returning to one, although the additional family still needs an untouched validation sample. Positive constant-background bounds survive expiry-specific exercise and mean allowances. In the discrete-outcome example, ATM-variance additivity can fail by much more than those allowances; good interpolation alone cannot establish structural meeting variances. Fitted second moments retain the constant-versus-flexible contrast with a 10\% extra moment allowance, but wider extrapolation bands remove the positive lower bounds. The two proxies expose different sources of uncertainty in the allocation.

A useful calibration therefore reports compatible meeting and group ranges alongside the restrictions used to obtain them. It tests withheld contracts and distinguishes a better fit from better predictions. The theoretical price-equivalence and factor-support results explain why some ambiguities persist even with richer pricing, while the numerical analysis shows where stronger modelling assumptions are doing the work.
